\documentclass[../../../include/open-logic-section]{subfiles}

\begin{document}

\olfileid[hi]{sth}{card-arithmetic}{opps}
\olsection{मूल संक्रियाओं की परिभाषा}

चूँकि हमें क्रम का ध्यान रखने की आवश्यकता नहीं है, कार्डिनल अंकगणित की
परिभाषा क्रमसूचक अंकगणित की अपेक्षा कुछ सरल है। हम योग, गुणन और घातांक को
एक साथ परिभाषित करेंगे।

\begin{defn}
जब $\cardfont{a}$ और $\cardfont{b}$ कार्डिनल हों:
\begin{align*}
	\cardfont{a} \cardplus \cardfont{b} &\defis 
	\card{\cardfont{a} \disjointsum \cardfont{b}}\\
	\cardfont{a} \cardtimes \cardfont{b} &\defis 
	\card{\cardfont{a} \times \cardfont{b}}\\
	\cardexpo{\cardfont{a}}{\cardfont{b}} &\defis 
	\card{\funfromto{\cardfont{b}}{\cardfont{a}}}
\end{align*}
जहाँ $\funfromto{X}{Y}=\Setabs{f}{f\text{ एक फलन है }X\to Y}$।
(यह दिखाना सरल है कि किन्हीं समुच्चयों $X$ और $Y$ के लिए
$\funfromto{X}{Y}$ का अस्तित्व है; इसे हम अभ्यास के लिए छोड़ते हैं।)
\end{defn}

\begin{prob}
$\Zminus$ में सिद्ध कीजिए कि किन्हीं समुच्चयों $X$ और $Y$ के लिए
$\funfromto{X}{Y}$ का अस्तित्व है। $\ZF$ में काम करते हुए
\olref[sth][ord-arithmetic][using-addition]{rankcomputation} की रीति से
$\setrank{X}$ और $\setrank{Y}$ से $\setrank{\funfromto{X}{Y}}$ परिकलित
कीजिए।
\end{prob}

इस परिभाषा की व्याख्या सहायक हो सकती है। योग के विषय में: इसमें
\olref[ord-arithmetic][add]{defdissum} में परिभाषित असंयुक्त योग
$\disjointsum$ की धारणा का उपयोग होता है; और यह देखना सरल है कि यह परिभाषा
परिमित स्थितियों में सही परिणाम देती है। गुणन के विषय में:
\olref[sfr][set][pai]{cardnmprod} हमें बताता है कि यदि $A$ के $n$ सदस्य और
$B$ के $m$ सदस्य हों, तो $A\times B$ के $n\cdot m$ सदस्य होते हैं, इसलिए
हमारी परिभाषा इस विचार का केवल परिमितेतर गुणन तक सामान्यीकरण करती है। घातांक
भी समान है: हम इस विचार का केवल परिमित से परिमितेतर तक सामान्यीकरण कर रहे
हैं। वास्तव में कुछ अर्थों में परिमितेतर कार्डिनल अंकगणित ``साधारण''
अंकगणित से परिमितेतर क्रमसूचक अंकगणित की अपेक्षा कहीं अधिक मिलता है:

\begin{prop}\ollabel{cardplustimescommute}
$\cardplus$ और $\cardtimes$ क्रमविनिमेय और साहचर्यात्मक हैं।
\end{prop}

\begin{proof}
क्रमविनिमेयता के लिए
\olref[cardinals][cardsasords]{lem:CardinalsBehaveRight} से यह देखना
पर्याप्त है कि $\cardeq{(\cardfont{a}\disjointsum
\cardfont{b})}{(\cardfont{b}\disjointsum\cardfont{a})}$ और
$\cardeq{(\cardfont{a}\times\cardfont{b})}{(\cardfont{b}\times
\cardfont{a})}$। साहचर्यता को हम अभ्यास के लिए छोड़ते हैं।
\end{proof}

\begin{prob}
सिद्ध कीजिए कि $\cardplus$ और $\cardtimes$ साहचर्यात्मक हैं।
\end{prob}

\begin{prop}
$A$ अनंत है तभी जब
$\card{A}\cardplus 1=1\cardplus\card{A}=\card{A}$।
\end{prop}

\begin{proof}
\olref[cardinals][classing]{generalinfinitycharacter} की तरह,
\olref[ord-arithmetic][using-addition]{ordinfinitycharacter} और
\olref[cardinals][cardsasords]{lem:CardinalsBehaveRight} से।
\end{proof}

इससे स्पष्ट होता है कि क्रमसूचक और कार्डिनल योग/गुणन के लिए हमें भिन्न
चिह्नों का उपयोग क्यों करना पड़ता है: ये वास्तव में \emph{भिन्न} संक्रियाएँ
हैं। अगले दो परिणाम दिखाते हैं कि क्रमसूचक और कार्डिनल घातांक भी भिन्न
संक्रियाएँ हैं। (याद करें कि \olref[z][infinity-again]{defnomega} से
$2=\{0,1\}$ निष्पन्न होता है):

\begin{lem}\ollabel{lem:SizePowerset2Exp}
किसी भी $A$ के लिए $\card{\Pow{A}}=\cardexpo{2}{\card{A}}$।
\end{lem}

\begin{proof}
प्रत्येक उपसमुच्चय $B\subseteq A$ के लिए
$\chi_B\in\funfromto{A}{2}$ इस प्रकार दिया हो:
\begin{align*}
	\chi_{B}(x) &\defis
	\begin{cases}
		1 & \text{यदि }x\in B\\
		0 & \text{अन्यथा।}
	\end{cases}
\end{align*}
अब $f(B)=\chi_B$ रखें; इससे कोई !!a{bijection}
$f\colon\Pow{A}\to\funfromto{A}{2}$ परिभाषित होता है। अतः
$\cardeq{\Pow{A}}{\funfromto{A}{2}}$। फलतः
$\cardeq{\Pow{A}}{\funfromto{\card{A}}{2}}$, इसलिए
$\card{\Pow{A}}=\card{\funfromto{\card{A}}{2}}=2^{\card{A}}$।
\end{proof}

यह संक्षिप्त प्रमाण वस्तुतः \olref[sfr][siz][red-alt]{sec} की चर्चा को
अपने में समेट लेता है। वहाँ हमने दिखाया था कि $\Pow{\omega}$ की अगणनीयता
को अनंत द्विआधारी शृंखलाओं के समुच्चय $\Bin^\omega$ की अगणनीयता तक कैसे
``घटाया'' जाए। वस्तुतः $\Bin^\omega$ केवल $\funfromto{\omega}{2}$ है; और
पिछले प्रमाण ने दिखाया कि \olref[sfr][siz][red-alt]{sec} में प्रयुक्त तर्क
$\omega$ के स्थान पर किसी भी समुच्चय $A$ के साथ चलेगा। यह परिणाम कार्डिनल
घातांक के बारे में एक त्वरित तथ्य भी देता है:

\begin{cor}\ollabel{cantorcor}
किसी भी कार्डिनल $\cardfont{a}$ के लिए
$\cardfont{a}<\cardexpo{2}{\cardfont{a}}$।
\end{cor}

\begin{proof}
कैंटर के प्रमेय (\olref[sfr][siz][car]{thm:cantor}) और
\olref{lem:SizePowerset2Exp} से।
\end{proof}
\noindent
अतः $\omega<\cardexpo{2}{\omega}$। किंतु ध्यान दें: यह \emph{कार्डिनल}
घातांक के बारे में परिणाम है। इसकी तुलना \emph{क्रमसूचक} घातांक से करनी
चाहिए, क्योंकि बाद वाली स्थिति में $\omega=\ordexpo{2}{\omega}$
(\olref[ord-arithmetic][expo]{sec} देखें)।

कार्डिनल घातांक की चर्चा करते हुए हम उस ``ढंग'' के बारे में भी कुछ अधिक
सटीक हो सकते हैं जिसमें $\Real$ !!{nonenumerable} है।

\begin{thm}\ollabel{continuumis2aleph0}
$\card{\Real} = \cardexpo{2}{\omega}$
\end{thm}

\begin{proof}[प्रमाण की रूपरेखा]
इसे सिद्ध करने के अनेक ढंग हैं। सबसे सीधा ढंग यह तर्क देना है कि
$\cardle{\Pow{\omega}}{\Real}$ और $\cardle{\Real}{\Pow{\omega}}$, फिर
श्रेडर--बर्नस्टीन से $\cardeq{\Real}{\Pow{\omega}}$ निष्कर्षित करना, और
\olref[card-arithmetic][opps]{lem:SizePowerset2Exp} से
$\card{\Real}=\cardexpo{2}{\omega}$ निष्कर्षित करना। अंतःक्षेप
$f\colon\Pow{\omega}\to\Real$ और
$g\colon\Real\to\Pow{\omega}$ परिभाषित करना हम एक (ज्ञानवर्धक) अभ्यास के
रूप में छोड़ते हैं।
\end{proof}

\begin{prob}
यह दिखाकर कि $\cardle{\Pow{\omega}}{\Real}$ और
$\cardle{\Real}{\Pow{\omega}}$, 
\olref[sth][card-arithmetic][opps]{continuumis2aleph0} का प्रमाण पूरा
कीजिए।
\end{prob}

\end{document}
